\chapter*{Exercice d'entraînement pour parcourir des listes}

\begin{Exercise}[title = {Le grand classique - Parcours par valeur}]
On dispose d'une liste notes contenant des nombres entiers. Écrire un script qui calcule et renvoie la moyenne de la liste.
\begin{itemize}
    \item \textbf{Contrainte :} Utiliser une boucle for element in notes:.
    \item la moyenne de la liste [10, 20, 15] doit afficher 15.0.
\end{itemize}

\end{Exercise}

\begin{Exercise}[title = {Recherche d'indice (Parcours par index)}]
    Écrire un script qui affiche l'\motcle{index} de la première occurrence d'une variable 'cible' dans une liste. Si l'élément n'est pas présent, le script doit afficher -1.
    \begin{itemize}
        \item \textbf{Contrainte :} Utilise for i in range(len(liste)): et une variable 'cible'
        \item \textbf{Exemple} : La liste [10, 50, 20] avec une cible de 50 doit afficher 1.
    \end{itemize}

\end{Exercise}

 \begin{Exercise}[title = {Transformation (Parcours par index)}]
     On veut doubler la valeur de chaque nombre dans une liste. Écrire un script  qui modifie une liste L.
     \begin{itemize}
         \item \textbf{Rappel :} Pour modifier une liste, on doit accéder aux cases via leurs indices : L[i] = ....
        \item Si ma\_liste = [1, 2, 3], le script doit afficher [2, 4, 6].
     \end{itemize}
 \end{Exercise}

\begin{Exercise}[title={Le maximum (Parcours par valeur)}]
    Sans utiliser la fonction max(), écrire un script qui parcoure la liste et renvoie le plus grand élément.
    \begin{itemize}
        \item \textbf{Indice :} Initialise une variable maximum avec le premier élément de la liste, puis compare-la avec les suivants.
    \end{itemize} 
\end{Exercise}

\newpage
\begin{Exercise}[title={Le Gestionnaire de Cinéma}]
On représente une petite salle de cinéma par un tableau à deux dimensions nommé salle.
\begin{itemize}
    \item Chaque sous-tableau est une rangée de sièges.
    \item Le chiffre 0 signifie que le siège est libre.
    \item Le chiffre 1 signifie que le siège est occupé.
\end{itemize}
\begin{lstlisting}
    salle = [
    [1, 0, 1, 1], # Rangée 0
    [0, 0, 1, 0], # Rangée 1
    [1, 1, 1, 1], # Rangée 2
    [0, 0, 0, 0]  # Rangée 3
    ]
\end{lstlisting}
\begin{itemize}
    \item \textbf{La mission :} Écrire un script qui parcourt une salle et renvoie le nombre total de sièges libres.
    \item Contraintes :
    \begin{itemize}
        \item On doit utiliser deux boucles imbriquées (une boucle for à l'intérieur d'une autre).
        \item La première boucle doit parcourir les rangées.
        \item La deuxième boucle doit parcourir les sièges de chaque rangée.
    \end{itemize}
\end{itemize}


\end{Exercise}

\section*{Les petits rappels}
\begin{itemize}
    \item \motcle{Parcours par valeur} (for x in L) : C'est le plus simple. A utiliser quand on a juste besoin de "lire" les données sans savoir où elles sont.
    \item \motcle{Parcours par index} (for i in range(len(L))) : Indispensable si on a besoin de l'\motcle{indice} (pour un affichage ou une recherche) ou si on doit \motcle{modifier} le contenu de la liste.
    \item Comment "entrer" dans le tableau ? Pour accéder à un élément précis, on utilise deux paires de crochets : salle[rangee][siege].
    \begin{itemize}
        \item salle[0] donne la première liste : [1, 0, 1, 1].
        \item salle[0][1] donne le deuxième élément de la première liste : 0.
    \end{itemize}
\end{itemize}
